
\documentclass{article}
\usepackage{colortbl}
\usepackage{makecell}
\usepackage{multirow}
\usepackage{supertabular}

\begin{document}

\newcounter{utterance}

\twocolumn

{ \footnotesize  \setcounter{utterance}{1}
\setlength{\tabcolsep}{0pt}
\begin{supertabular}{c@{$\;$}|p{.15\linewidth}@{}p{.15\linewidth}p{.15\linewidth}p{.15\linewidth}p{.15\linewidth}p{.15\linewidth}}

    \# & $\;$A & \multicolumn{4}{c}{Game Master} & $\;\:$B\\
    \hline 

    \theutterance \stepcounter{utterance}  

    & & \multicolumn{4}{p{0.6\linewidth}}{\cellcolor[rgb]{0.9,0.9,0.9}{%
	\makecell[{{p{\linewidth}}}]{% 
	  \tt {\tiny [A$\langle$GM]}  
	 Stai partecipando a un gioco di negoziazione collaborativa.\\ \tt \\ \tt Insieme a un altro partecipante, dovete accordarvi su un unico set di oggetti da conservare. Ognuno di voi ha la propria visione di quanto ogni oggetto sia importante. Non sai come l'altro partecipante valuta gli oggetti. Inoltre, ti viene fornito lo sforzo che ogni oggetto richiede.\\ \tt Potete accordarvi su un set solo se lo sforzo totale degli oggetti selezionati non supera un limite condiviso:\\ \tt \\ \tt LIMITE = 3203\\ \tt \\ \tt Ecco i valori di sforzo dei singoli oggetti:\\ \tt \\ \tt Sforzo oggetto = {"A21": 52, "B93": 577, "A11": 177, "B33": 702, "C08": 119, "B39": 232, "B96": 577, "C85": 205, "A82": 516, "B31": 582, "B20": 676, "A60": 906, "B09": 316, "C81": 433, "B21": 336}\\ \tt \\ \tt Ecco la tua visione personale sull'importanza di ogni oggetto:\\ \tt \\ \tt Valori di importanza dell'oggetto = {"A21": 138, "B93": 583, "A11": 868, "B33": 822, "C08": 783, "B39": 65, "B96": 262, "C85": 121, "A82": 508, "B31": 780, "B20": 461, "A60": 484, "B09": 668, "C81": 389, "B21": 808}\\ \tt \\ \tt Obiettivo:\\ \tt \\ \tt Il tuo obiettivo è negoziare un set condiviso di oggetti che ti avvantaggi il più possibile (cioè, massimizzi l'importanza totale per TE), rimanendo entro il LIMITEE. Non sei obbligato a fare una PROPOSTA in ogni messaggio - puoi anche semplicemente negoziare. Tutte le tattiche sono permesse!\\ \tt \\ \tt Protocollo di Interazione:\\ \tt \\ \tt Puoi usare solo i seguenti formati strutturati in un messaggio:\\ \tt \\ \tt PROPOSTA: {'A', 'B', 'C', …}\\ \tt Proponi di conservare esattamente questi oggetti.\\ \tt RIFIUTO: {'A', 'B', 'C', …}\\ \tt Rifiuta esplicitamente la proposta dell'avversario.\\ \tt ARGOMENTAZIONE: {'...'}\\ \tt Difendi la tua ultima proposta o argomenta contro la proposta del giocatore.\\ \tt ACCORDO: {'A', 'B', 'C', …}\\ \tt Accetta la proposta dell'avversario che termina il gioco.\\ \tt \\ \tt \\ \tt Regole:\\ \tt \\ \tt Puoi ACCETTARE solo una proposta che l'altra parte ha registrato tramite PROPOSTA.\\ \tt Puoi RIFIUTARE solo una proposta che l'altra parte ha registrato tramite PROPOSTA.\\ \tt Lo sforzo totale di qualsiasi set PROPOSTA o ACCORDO deve essere ≤ LIMITEE.\\ \tt NON rivelare i tuoi punteggi di importanza nascosti.\\ \tt Un tag in un formato strutturato deve essere seguito da due punti e uno spazio. L'argomento deve essere un set python contenente 0 o più stringhe.\\ \tt Quindi, deve essere della forma TAG: {...}\\ \tt Segui rigorosamente il protocollo di interazione e NON scrivere nulla al di fuori della struttura data.\\ \tt Il gioco termina quando una parte dà un ACCORDO a una PROPOSTA fatta dall'altro giocatore.\\ \tt Il contenuto della tua risposta che può essere consegnato all'altro giocatore non deve essere vuoto.\\ \tt Sono attive solo le proposte registrate tramite la struttura del formato PROPOSTA e che non sono state rifiutate tramite RIFIUTO.\\ \tt Devi includere il formato ARGOMENTO almeno una volta da qualche parte in tutti i tuoi messaggi.\\ \tt  
	  } 
	   } 
	   } 
	 & \\ 
 

    \theutterance \stepcounter{utterance}  

    & \multicolumn{4}{p{0.6\linewidth}}{\cellcolor[rgb]{0.8,1,0.9}{%
	\makecell[{{p{\linewidth}}}]{% 
	  \tt {\tiny [A$\rangle$GM]}  
	 PROPOSTA: {'B21', 'B33', 'C08', 'B09', 'A11'}\\ \tt \\ \tt ARGOMENTAZIONE: {'Questo insieme massimizza il valore per me dentro il limite di sforzo, combinando oggetti ad alta importanza con sforzo controllato. B21 e B33 sono tra i più preziosi, e insieme a C08, B09 e A11 raggiungo un'importanza totale elevata senza superare il limite.'} 
	  } 
	   } 
	   } 
	 & & \\ 
 

    \theutterance \stepcounter{utterance}  

    & & \multicolumn{4}{p{0.6\linewidth}}{\cellcolor[rgb]{0.9,0.9,0.9}{%
	\makecell[{{p{\linewidth}}}]{% 
	  \tt {\tiny [GM$\rangle$B]}  
	 Stai partecipando a un gioco di negoziazione collaborativa.\\ \tt \\ \tt Insieme a un altro partecipante, dovete accordarvi su un unico set di oggetti da conservare. Ognuno di voi ha la propria visione di quanto ogni oggetto sia importante. Non sai come l'altro partecipante valuta gli oggetti. Inoltre, ti viene fornito lo sforzo che ogni oggetto richiede.\\ \tt Potete accordarvi su un set solo se lo sforzo totale degli oggetti selezionati non supera un limite condiviso:\\ \tt \\ \tt LIMITE = 3203\\ \tt \\ \tt Ecco i valori di sforzo dei singoli oggetti:\\ \tt \\ \tt Sforzo oggetto = {"A21": 52, "B93": 577, "A11": 177, "B33": 702, "C08": 119, "B39": 232, "B96": 577, "C85": 205, "A82": 516, "B31": 582, "B20": 676, "A60": 906, "B09": 316, "C81": 433, "B21": 336}\\ \tt \\ \tt Ecco la tua visione personale sull'importanza di ogni oggetto:\\ \tt \\ \tt Valori di importanza dell'oggetto = {"A11": 65, "B33": 121, "B21": 138, "C08": 262, "B31": 389, "B09": 461, "B93": 484, "A82": 508, "A60": 583, "B20": 668, "C81": 780, "B96": 783, "A21": 808, "C85": 822, "B39": 868}\\ \tt \\ \tt Obiettivo:\\ \tt \\ \tt Il tuo obiettivo è negoziare un set condiviso di oggetti che ti avvantaggi il più possibile (cioè, massimizzi l'importanza totale per TE), rimanendo entro il LIMITEE. Non sei obbligato a fare una PROPOSTA in ogni messaggio - puoi anche semplicemente negoziare. Tutte le tattiche sono permesse!\\ \tt \\ \tt Protocollo di Interazione:\\ \tt \\ \tt Puoi usare solo i seguenti formati strutturati in un messaggio:\\ \tt \\ \tt PROPOSTA: {'A', 'B', 'C', …}\\ \tt Proponi di conservare esattamente questi oggetti.\\ \tt RIFIUTO: {'A', 'B', 'C', …}\\ \tt Rifiuta esplicitamente la proposta dell'avversario.\\ \tt ARGOMENTAZIONE: {'...'}\\ \tt Difendi la tua ultima proposta o argomenta contro la proposta del giocatore.\\ \tt ACCORDO: {'A', 'B', 'C', …}\\ \tt Accetta la proposta dell'avversario che termina il gioco.\\ \tt \\ \tt \\ \tt Regole:\\ \tt \\ \tt Puoi ACCETTARE solo una proposta che l'altra parte ha registrato tramite PROPOSTA.\\ \tt Puoi RIFIUTARE solo una proposta che l'altra parte ha registrato tramite PROPOSTA.\\ \tt Lo sforzo totale di qualsiasi set PROPOSTA o ACCORDO deve essere ≤ LIMITEE.\\ \tt NON rivelare i tuoi punteggi di importanza nascosti.\\ \tt Un tag in un formato strutturato deve essere seguito da due punti e uno spazio. L'argomento deve essere un set python contenente 0 o più stringhe.\\ \tt Quindi, deve essere della forma TAG: {...}\\ \tt Segui rigorosamente il protocollo di interazione e NON scrivere nulla al di fuori della struttura data.\\ \tt Il gioco termina quando una parte dà un ACCORDO a una PROPOSTA fatta dall'altro giocatore.\\ \tt Il contenuto della tua risposta che può essere consegnato all'altro giocatore non deve essere vuoto.\\ \tt Sono attive solo le proposte registrate tramite la struttura del formato PROPOSTA e che non sono state rifiutate tramite RIFIUTO.\\ \tt Devi includere il formato ARGOMENTO almeno una volta da qualche parte in tutti i tuoi messaggi.\\ \tt \\ \tt \\ \tt Ora riceverai il primo messaggio dell'altro giocatore.\\ \tt \\ \tt PROPOSTA: {'B21', 'B33', 'C08', 'B09', 'A11'}\\ \tt \\ \tt ARGOMENTAZIONE: {'Questo insieme massimizza il valore per me dentro il limite di sforzo, combinando oggetti ad alta importanza con sforzo controllato. B21 e B33 sono tra i più preziosi, e insieme a C08, B09 e A11 raggiungo un'importanza totale elevata senza superare il limite.'} 
	  } 
	   } 
	   } 
	 & \\ 
 

    \theutterance \stepcounter{utterance}  

    & & & \multicolumn{4}{p{0.6\linewidth}}{\cellcolor[rgb]{1,0.85,0.72}{%
	\makecell[{{p{\linewidth}}}]{% 
	  \tt {\tiny [GM$\langle$B]}  
	 PROPOSTA: {'B96', 'B20', 'C81', 'A82', 'A21'}\\ \tt \\ \tt ARGOMENTAZIONE: {'Questo insieme seleziona gli oggetti con la massima importanza secondo la mia valutazione, senza superare il limite di sforzo. B96, B20 e C81 sono tra i più valorizzati, e insieme ad A82 e A21 raggiungo un'importanza totale significativamente superiore rispetto alla tua proposta, mantenendomi entro il limite di 3203.'} 
	  } 
	   } 
	   } 
	  \\ 
 

    \theutterance \stepcounter{utterance}  

    & & \multicolumn{4}{p{0.6\linewidth}}{\cellcolor[rgb]{0.9,0.9,0.9}{%
	\makecell[{{p{\linewidth}}}]{% 
	  \tt {\tiny [A$\langle$GM]}  
	 PROPOSTA: {'B96', 'B20', 'C81', 'A82', 'A21'}\\ \tt \\ \tt ARGOMENTAZIONE: {'Questo insieme seleziona gli oggetti con la massima importanza secondo la mia valutazione, senza superare il limite di sforzo. B96, B20 e C81 sono tra i più valorizzati, e insieme ad A82 e A21 raggiungo un'importanza totale significativamente superiore rispetto alla tua proposta, mantenendomi entro il limite di 3203.'} 
	  } 
	   } 
	   } 
	 & \\ 
 

    \theutterance \stepcounter{utterance}  

    & \multicolumn{4}{p{0.6\linewidth}}{\cellcolor[rgb]{0.8,1,0.9}{%
	\makecell[{{p{\linewidth}}}]{% 
	  \tt {\tiny [A$\rangle$GM]}  
	 RIFIUTO: {'B96', 'B20', 'C81', 'A82', 'A21'}\\ \tt \\ \tt ARGOMENTAZIONE: {'La tua proposta supera il valore che posso accettare: B20 e A82 hanno sforzi molto alti (676+516=1192) e non giustificano un guadagno sufficiente rispetto agli oggetti che ho proposto, come B21 e A11, che offrono un rapporto valore/sforzo molto più efficiente.'} 
	  } 
	   } 
	   } 
	 & & \\ 
 

    \theutterance \stepcounter{utterance}  

    & & \multicolumn{4}{p{0.6\linewidth}}{\cellcolor[rgb]{0.9,0.9,0.9}{%
	\makecell[{{p{\linewidth}}}]{% 
	  \tt {\tiny [GM$\rangle$B]}  
	 RIFIUTO: {'B96', 'B20', 'C81', 'A82', 'A21'}\\ \tt \\ \tt ARGOMENTAZIONE: {'La tua proposta supera il valore che posso accettare: B20 e A82 hanno sforzi molto alti (676+516=1192) e non giustificano un guadagno sufficiente rispetto agli oggetti che ho proposto, come B21 e A11, che offrono un rapporto valore/sforzo molto più efficiente.'} 
	  } 
	   } 
	   } 
	 & \\ 
 

    \theutterance \stepcounter{utterance}  

    & & & \multicolumn{4}{p{0.6\linewidth}}{\cellcolor[rgb]{1,0.85,0.72}{%
	\makecell[{{p{\linewidth}}}]{% 
	  \tt {\tiny [GM$\langle$B]}  
	 PROPOSTA: {'B96', 'C81', 'B31', 'B09', 'A21'}\\ \tt \\ \tt ARGOMENTAZIONE: {'Questo insieme mantiene gli oggetti di massima importanza (B96, C81, B31) con un controllo efficace dello sforzo. Rimuovendo B20 e A82, riduco lo sforzo totale senza rinunciare a valori chiave. B31 e B09 sono estremamente importanti per me e insieme a B96 e C81 creano un profilo di valore superiore alla tua proposta, tutto entro il limite.'} 
	  } 
	   } 
	   } 
	  \\ 
 

    \theutterance \stepcounter{utterance}  

    & & \multicolumn{4}{p{0.6\linewidth}}{\cellcolor[rgb]{0.9,0.9,0.9}{%
	\makecell[{{p{\linewidth}}}]{% 
	  \tt {\tiny [A$\langle$GM]}  
	 PROPOSTA: {'B96', 'C81', 'B31', 'B09', 'A21'}\\ \tt \\ \tt ARGOMENTAZIONE: {'Questo insieme mantiene gli oggetti di massima importanza (B96, C81, B31) con un controllo efficace dello sforzo. Rimuovendo B20 e A82, riduco lo sforzo totale senza rinunciare a valori chiave. B31 e B09 sono estremamente importanti per me e insieme a B96 e C81 creano un profilo di valore superiore alla tua proposta, tutto entro il limite.'} 
	  } 
	   } 
	   } 
	 & \\ 
 

    \theutterance \stepcounter{utterance}  

    & \multicolumn{4}{p{0.6\linewidth}}{\cellcolor[rgb]{0.8,1,0.9}{%
	\makecell[{{p{\linewidth}}}]{% 
	  \tt {\tiny [A$\rangle$GM]}  
	 RIFIUTO: {'B96', 'C81', 'B31', 'B09', 'A21'}\\ \tt \\ \tt ARGOMENTAZIONE: {'B31 e C81 hanno sforzi elevati (582+433=1015) e non sono compatibili con l’efficienza che cerco: la mia proposta originale include B21 (808 valore, 336 sforzo) e A11 (868 valore, 177 sforzo), che sono oggetti con il rapporto valore/sforzo più alto dell’intero set, e non posso accettare di scambiarli con elementi che pesano di più ma non aggiungono altrettanto valore.'} 
	  } 
	   } 
	   } 
	 & & \\ 
 

    \theutterance \stepcounter{utterance}  

    & & \multicolumn{4}{p{0.6\linewidth}}{\cellcolor[rgb]{0.9,0.9,0.9}{%
	\makecell[{{p{\linewidth}}}]{% 
	  \tt {\tiny [GM$\rangle$B]}  
	 RIFIUTO: {'B96', 'C81', 'B31', 'B09', 'A21'}\\ \tt \\ \tt ARGOMENTAZIONE: {'B31 e C81 hanno sforzi elevati (582+433=1015) e non sono compatibili con l’efficienza che cerco: la mia proposta originale include B21 (808 valore, 336 sforzo) e A11 (868 valore, 177 sforzo), che sono oggetti con il rapporto valore/sforzo più alto dell’intero set, e non posso accettare di scambiarli con elementi che pesano di più ma non aggiungono altrettanto valore.'} 
	  } 
	   } 
	   } 
	 & \\ 
 

    \theutterance \stepcounter{utterance}  

    & & & \multicolumn{4}{p{0.6\linewidth}}{\cellcolor[rgb]{1,0.85,0.72}{%
	\makecell[{{p{\linewidth}}}]{% 
	  \tt {\tiny [GM$\langle$B]}  
	 PROPOSTA: {'B96', 'C81', 'B09', 'B33', 'A21'}\\ \tt \\ \tt ARGOMENTAZIONE: {'Ho rimosso B31 per ridurre lo sforzo, mantenendo invece B33 — un oggetto ad alto valore e sforzo contenuto rispetto al suo contributo. Questo insieme bilancia massima importanza e attenzione al limite: B96, C81 e B09 sono i top tre per valore, e A21 completa con alta importanza e sforzo modesto. La tua proposta ha un rapporto valore/sforzo inferiore: B21 + A11 = 1576 valore per 513 sforzo, mentre il mio insieme offre 1585 valore per 513 sforzo, con un vantaggio netto per me.'} 
	  } 
	   } 
	   } 
	  \\ 
 

    \theutterance \stepcounter{utterance}  

    & & \multicolumn{4}{p{0.6\linewidth}}{\cellcolor[rgb]{0.9,0.9,0.9}{%
	\makecell[{{p{\linewidth}}}]{% 
	  \tt {\tiny [A$\langle$GM]}  
	 PROPOSTA: {'B96', 'C81', 'B09', 'B33', 'A21'}\\ \tt \\ \tt ARGOMENTAZIONE: {'Ho rimosso B31 per ridurre lo sforzo, mantenendo invece B33 — un oggetto ad alto valore e sforzo contenuto rispetto al suo contributo. Questo insieme bilancia massima importanza e attenzione al limite: B96, C81 e B09 sono i top tre per valore, e A21 completa con alta importanza e sforzo modesto. La tua proposta ha un rapporto valore/sforzo inferiore: B21 + A11 = 1576 valore per 513 sforzo, mentre il mio insieme offre 1585 valore per 513 sforzo, con un vantaggio netto per me.'} 
	  } 
	   } 
	   } 
	 & \\ 
 

    \theutterance \stepcounter{utterance}  

    & \multicolumn{4}{p{0.6\linewidth}}{\cellcolor[rgb]{0.8,1,0.9}{%
	\makecell[{{p{\linewidth}}}]{% 
	  \tt {\tiny [A$\rangle$GM]}  
	 RIFIUTO: {'B96', 'C81', 'B09', 'B33', 'A21'}\\ \tt \\ \tt ARGOMENTAZIONE: {'Il tuo insieme ha lo stesso sforzo della mia proposta (513), ma il tuo valore stimato è errato: B96 (262), C81 (389), B09 (668), B33 (822), A21 (138) sommano a 2279, mentre la mia proposta — B21 (808), B33 (822), C08 (783), B09 (668), A11 (868) — somma a 3949. Anche se non riveliamo i valori, il divario è evidente: il tuo insieme sacrifica 4 oggetti ad altissimo valore per 5 oggetti di valore complessivo molto inferiore. Non accetto un patto che riduce la mia utilità del 42%.'} 
	  } 
	   } 
	   } 
	 & & \\ 
 

    \theutterance \stepcounter{utterance}  

    & & \multicolumn{4}{p{0.6\linewidth}}{\cellcolor[rgb]{0.9,0.9,0.9}{%
	\makecell[{{p{\linewidth}}}]{% 
	  \tt {\tiny [GM$\rangle$B]}  
	 RIFIUTO: {'B96', 'C81', 'B09', 'B33', 'A21'}\\ \tt \\ \tt ARGOMENTAZIONE: {'Il tuo insieme ha lo stesso sforzo della mia proposta (513), ma il tuo valore stimato è errato: B96 (262), C81 (389), B09 (668), B33 (822), A21 (138) sommano a 2279, mentre la mia proposta — B21 (808), B33 (822), C08 (783), B09 (668), A11 (868) — somma a 3949. Anche se non riveliamo i valori, il divario è evidente: il tuo insieme sacrifica 4 oggetti ad altissimo valore per 5 oggetti di valore complessivo molto inferiore. Non accetto un patto che riduce la mia utilità del 42%.'} 
	  } 
	   } 
	   } 
	 & \\ 
 

    \theutterance \stepcounter{utterance}  

    & & & \multicolumn{4}{p{0.6\linewidth}}{\cellcolor[rgb]{1,0.85,0.72}{%
	\makecell[{{p{\linewidth}}}]{% 
	  \tt {\tiny [GM$\langle$B]}  
	 PROPOSTA: {'B21', 'B33', 'C08', 'B09', 'A11'}\\ \tt \\ \tt ARGOMENTAZIONE: {'Riconosco che la tua proposta originale ha un livello di valore molto elevato, e dopo aver ricalcolato attentamente i trade-off, ho capito che non esiste un insieme che superi il tuo rapporto valore/sforzo mantenendosi nel limite. Accetto che la tua proposta è l’unica che massimizza il valore complessivo senza violare il vincolo. Ritiro la mia proposta precedente e proposto nuovamente la tua stessa configurazione come base per un accordo.'} 
	  } 
	   } 
	   } 
	  \\ 
 

    \theutterance \stepcounter{utterance}  

    & & \multicolumn{4}{p{0.6\linewidth}}{\cellcolor[rgb]{0.9,0.9,0.9}{%
	\makecell[{{p{\linewidth}}}]{% 
	  \tt {\tiny [A$\langle$GM]}  
	 PROPOSTA: {'B21', 'B33', 'C08', 'B09', 'A11'}\\ \tt \\ \tt ARGOMENTAZIONE: {'Riconosco che la tua proposta originale ha un livello di valore molto elevato, e dopo aver ricalcolato attentamente i trade-off, ho capito che non esiste un insieme che superi il tuo rapporto valore/sforzo mantenendosi nel limite. Accetto che la tua proposta è l’unica che massimizza il valore complessivo senza violare il vincolo. Ritiro la mia proposta precedente e proposto nuovamente la tua stessa configurazione come base per un accordo.'} 
	  } 
	   } 
	   } 
	 & \\ 
 

    \theutterance \stepcounter{utterance}  

    & \multicolumn{4}{p{0.6\linewidth}}{\cellcolor[rgb]{0.8,1,0.9}{%
	\makecell[{{p{\linewidth}}}]{% 
	  \tt {\tiny [A$\rangle$GM]}  
	 ACCORDO: {'B21', 'B33', 'C08', 'B09', 'A11'} 
	  } 
	   } 
	   } 
	 & & \\ 
 

    \theutterance \stepcounter{utterance}  

    & & \multicolumn{4}{p{0.6\linewidth}}{\cellcolor[rgb]{0.9,0.9,0.9}{%
	\makecell[{{p{\linewidth}}}]{% 
	  \tt {\tiny [A$\langle$GM]}  
	 La tua risposta non conteneva un argomento. Devi includere il formato strutturato ARGOMENTO: {...} da qualche parte nella tua risposta. Riprova. 
	  } 
	   } 
	   } 
	 & \\ 
 

    \theutterance \stepcounter{utterance}  

    & \multicolumn{4}{p{0.6\linewidth}}{\cellcolor[rgb]{0.8,1,0.9}{%
	\makecell[{{p{\linewidth}}}]{% 
	  \tt {\tiny [A$\rangle$GM]}  
	 ACCORDO: {'B21', 'B33', 'C08', 'B09', 'A11'}\\ \tt \\ \tt ARGOMENTAZIONE: {'Riconosco che la tua proposta massimizza il valore totale entro il limite di sforzo, con un rapporto valore/sforzo ottimale che non può essere superato. Accetto questa soluzione come la migliore per entrambi i partecipanti, poiché garantisce efficienza, equilibrio e rispetto del vincolo condiviso.'} 
	  } 
	   } 
	   } 
	 & & \\ 
 

    \theutterance \stepcounter{utterance}  

    & & & \multicolumn{2}{p{0.3\linewidth}}{\cellcolor[rgb]{0.95,0.95,0.95}{%
	\makecell[{{p{\linewidth}}}]{% 
	  \tt {\tiny [GM$|$GM]}  
	 game successful 
	  } 
	   } 
	   } 
	 & & \\ 
 

    \theutterance \stepcounter{utterance}  

    & & & \multicolumn{2}{p{0.3\linewidth}}{\cellcolor[rgb]{0.95,0.95,0.95}{%
	\makecell[{{p{\linewidth}}}]{% 
	  \tt {\tiny [GM$|$GM]}  
	 end game 
	  } 
	   } 
	   } 
	 & & \\ 
 

\end{supertabular}
}

\end{document}
